% 中文关键字
\newcommand{\keywordsCn}[1][遥感图像；微小船舶检测；LEVIR-Ship；目标检测；检测系统]{#1}
% 中文摘要
\begin{abstract}
    遥感图像中的微小船舶检测常服务于海上监管、港口监测、海事执法和应急救援等任务。与普通自然图像不同，中分辨率遥感图像里船舶目标占像素少、纹理弱，附近又可能出现海浪、云层、岸线和港区设施等干扰，模型容易把背景当作船，也容易漏掉弱小目标。围绕这一类场景，本文以 LEVIR-Ship 数据集为实验对象，对不同检测路线进行比较，并在实验结果基础上实现一套面向毕业设计场景的检测系统。

    算法部分先整理通用目标检测、小目标检测和遥感船舶检测的相关研究，再固定数据划分、评价指标和结果记录方式，选取 DRENet、FCOS 与 YOLO26 进行实验。本文以 AP50 作为主指标，同时记录 AP50-95、Precision、Recall 和 F1。测试结果中，DRENet 的 AP50 为 0.7949，Recall 为 0.8511，说明它更愿意保留候选目标；YOLO26 的 AP50 为 0.7950，AP50-95 为 0.3170，Precision 为 0.8430，误检控制和系统接入表现更稳；FCOS 的 AP50 为 0.7700，主要用于保留 anchor-free 路线的参照。阈值消融和输入尺寸实验进一步显示，三类模型在漏检、误检和尺度扰动下的反应并不一致。

    系统部分将本地推理能力封装为轻量级遥感微小船舶检测系统。系统前端采用 React，负责图像上传、任务查询和结果展示；后端采用 FastAPI，负责推理服务和任务接口；SQLite 保存模型信息、任务状态、检测框结果和输出文件路径。系统已经支持单图与批量任务、模型启停、异步执行、可视化结果回看以及历史任务查询，可以把离线实验中的检测结果放到可操作的页面流程中复查。

    通过上述实验和系统实现，本文得到的判断是：在遥感微小船舶任务中，单看一个精度数值不足以解释模型行为，必须同时观察阈值、输入尺寸、误检样例和系统运行口径。本文保留下来的实验记录、分析过程和系统原型，可作为后续融合评测、模型优化与部署整理的基础。
\end{abstract}

% 英文关键字
\newcommand{\keywordsEn}[1][remote sensing images, tiny ship detection, LEVIR-Ship, object detection, detection system]{#1}
% 英文摘要
\begin{abstractEn}
    Tiny ship detection in remote sensing imagery is a basic component of maritime surveillance, port monitoring, law enforcement and emergency response. In mid-resolution images, ships often occupy only a few pixels and have weak texture cues, while waves, clouds, coastlines and harbor facilities introduce strong background interference. These factors make missed detections, false alarms and unstable cross-scene performance frequent in practice. Based on this setting, this thesis conducts comparative experiments on the public LEVIR-Ship dataset and implements a lightweight detection system for the project scenario.

    For the algorithmic study, this thesis reviews related work on general object detection, small object detection and remote-sensing ship detection, and then builds a unified protocol for data split, evaluation metrics and result recording. Under this protocol, DRENet, FCOS and YOLO26 are selected for comparison. AP50 is used as the primary metric, with AP50-95, Precision, Recall and F1 used as complementary indicators. The results show that DRENet achieves an AP50 of 0.7949 and a Recall of 0.8511 on LEVIR-Ship, indicating strong target discovery ability. YOLO26 achieves an AP50 of 0.7950, an AP50-95 of 0.3170 and a Precision of 0.8430, giving a relatively balanced result among accuracy, localization quality and deployability. FCOS achieves an AP50 of 0.7700 and is kept as an anchor-free reference baseline. Threshold ablation and input-size sensitivity analysis further show that the three paradigms differ in false-alarm suppression, missed-detection control and scale adaptability.

    For the system part, the existing inference capability is wrapped into a lightweight remote-sensing tiny ship detection system. The frontend is built with React for task submission and result visualization, the backend uses FastAPI for inference services and task management, and SQLite stores model information, task status, detection results and artifact paths. The system supports single-image and batch-image submission, model activation management, asynchronous task execution, detection-box visualization, structured result review and history query, so that offline inference outputs can be inspected through an interactive workflow.

    The study suggests that unified multi-model evaluation, threshold ablation, input-size sensitivity analysis and system implementation together provide a clearer view of the applicable boundaries of different detection paradigms in remote-sensing tiny ship detection. The experimental records, analysis process and system prototype in this thesis can support later work on model optimization, fusion evaluation and engineering deployment.
\end{abstractEn}
